\paragraph{Anticipación}

La objeción más creíble al diseño es la anticipación. La construcción de una
estación de servicio es visible durante meses, de modo que una incumbente podría
comenzar a ajustar sus precios antes de la apertura. Si así fuera, el salto
estimado sería solo la cola de una respuesta iniciada antes, y los bins de seis
meses de la especificación principal no permitirían verlo.

La \Cref{fig:anticipacion} lleva el estudio de eventos a resolución semanal en
torno a la semana exacta de apertura de la entrante, con veintiséis semanas a
cada lado. Donde la figura principal tiene un solo bin previo, esta tiene
veinticinco coeficientes.

\begin{figure}[H]
\centering
\includegraphics[width=\textwidth]{anticipacion.pdf}
\caption{Trayectoria semanal en torno a la apertura de la entrante}
\label{fig:anticipacion}
\end{figure}

No hay movimiento previo. El test conjunto sobre las veinticinco semanas
anteriores no rechaza —$p = 0{,}85$ en la gasolina de 93 octanos y $p = 0{,}57$
en el diésel— y la pendiente de la trayectoria previa, que es el contraste con
potencia frente a una deriva suave, es de seis diezmilésimas de punto porcentual
por semana en la 93 y de tres milésimas en el diésel, ninguna distinguible de
cero. Conviene notar, además, que la anticipación es una amenaza de dirección
conocida: de existir, los precios ya estarían más bajos en el período de
referencia y el efecto estimado sería una cota inferior en magnitud.

\paragraph{Predictibilidad del momento de entrada}

La verificación más directa del supuesto consiste en preguntar si las
condiciones locales predicen \emph{cuándo} llega la entrada. Se estima un modelo
de duración en tiempo discreto sobre el conjunto de incumbentes en riesgo, con
el margen del mercado local, su variación sobre el año previo, el nivel de
precios local y el número de competidores, todos medidos con rezago. El
ejercicio se hace sobre dos muestras: todos los incumbentes en riesgo, donde
las condiciones locales sí deberían predecir la entrada porque ahí opera la
selección de ubicación; y solo aquellos que efectivamente reciben una entrada,
donde el supuesto exige que no la predigan.

El resultado separa con nitidez ambos canales. Las variables de rentabilidad
—margen local, su variación y el nivel de precios— no predicen ni la incidencia
ni el momento: el test conjunto arroja $p = 0{,}96$ en la primera muestra y
$p = 0{,}44$ en la segunda. Lo único que predice la entrada es el número de
competidores, en ambas muestras. Es decir, las entrantes se orientan por la
densidad del mercado, tal como anticipa la \Cref{prop:libre}, y no por el margen
que las incumbentes venían capturando, que es el canal que sesgaría la
estimación de precios. Queda como limitación reconocida que la densificación de
un mercado podría venir acompañada de crecimiento de la demanda local, algo que
los datos disponibles no permiten descartar; la sección de limitaciones vuelve
sobre este punto.
